Maintenance rules often trigger work when a lower prediction bound on remaining useful life (RUL) falls to a lead time, yet interval quality is judged by clean-input coverage. This stress test asks how structured sensor degradation changes coverage and trigger errors. On C-MAPSS (four sub-datasets, two backbones, five interval mechanisms), noise, bias, drift and gain errors were applied to the sensor channels of each test engine's final window, and a one-sided 95\% lower bound was compared with lead times of 10, 20 and 30 cycles. At matched out-of-range input share, coverage differed widely between perturbation families: on FD002 at 10\% out-of-range inputs, drift gave \JMatchDrift{} and bias \JMatchBias. A positive drift on all sensor channels caused premature triggers in \JDriftPremRange\% of evaluated six-condition engines at a 20-cycle lead time and no unrecognised need, whereas a constant negative offset caused unrecognised needs in \JCohNegUnrecSix\%. Output swaps identified the mean shift as the main driver, with scale changes amplifying or offsetting it. Cost rankings of mechanisms changed across cost ratios in \JRevTotal{} of \JRevGroups{} settings, never under drift. Results are terminal-instant, simulated and partly exploratory; they support reporting both trigger errors beside coverage, not deployment guarantees.
